\setcounter{section}{5}

\Needspace{5\baselineskip}
\hypertarget{the-muon-optimizer}{%
\section{The Muon Optimizer}\label{the-muon-optimizer}}

Adam/AdamW implements adaptive learning rates by scaling each dimension (parameter) separately through division by a scalar. In optimization, however, these dimensions represent only one coordinate system in a high-dimensional parameter space. Muon updates parameters by orthogonalizing the momentum matrix, making update matrix \(U\) approximately orthogonal, satisfying \(U^{T}U\approx I\).

Geometrically, it normalizes along all principal directions (singular-vector directions) in parameter space, taking equal-length steps along them. This resembles Newton's method performing gradient descent with equal step lengths in each direction of a transformed space (although these singular-vector directions are not Newton's directions), rather than Adam's attempt merely to adjust movement along individual coordinate axes.

For matrix orthogonalization, Muon uses Newton--Schulz iteration to approximate the matrix ``sign function,'' greatly reducing computation.

Some large-model training now uses modified Muon variants. Kimi K2, for example, uses MuonClip, which adds QK-clip to Muon, rather than the unmodified original optimizer.

\FloatBarrier
\Needspace{5\baselineskip}
\hypertarget{references-30}{%
\subsection{References}\label{references-30}}

\vspace{0.25em}
\begin{itemize}
\tightlist
\item
  Jordan, K., et al.~(2024). \href{https://kellerjordan.github.io/posts/muon/}{Muon: An optimizer for hidden layers in neural networks}.
\item
  Kimi Team. (2025). \href{https://arxiv.org/abs/2507.20534}{Kimi K2: Open Agentic Intelligence}. arXiv:2507.20534.
\end{itemize}

\clearpage
